% ID: SEL-D0017-BC-E06
% ESTADO: cotejada. Segunda pasada visual completa y borrador compilado revisados.
% Fuente: Selectividad/matematicas_aplicadas/matematicas_aplicadas_P_E.pdf
% SHA-256: bc88bfeb0f184eb0c806c4efe04e84c24f3e5fcef9caaa12e94fd48d9f67a389
% Página: 2; bloque C; ejercicio original 6.

El censo de una población andaluza está compuesto en total por $15\,000$
personas, de las cuales $8\,500$ son mujeres. Se sabe que el 15\% de las
mujeres y el 20\% de los hombres censados en dicha población han viajado
alguna vez a un país extranjero. Se elige al azar una persona censada en
dicha población.
\begin{enumerate}[label=\alph*)]
  \item \textbf{(1.25 puntos)} \textquestiondown{}Cuál es la probabilidad de
  que haya viajado al extranjero?
  \item \textbf{(1.25 puntos)} Si se sabe que esta persona no ha viajado
  al extranjero, \textquestiondown{}cuál es la probabilidad de que sea hombre?
\end{enumerate}
